\section{Discussion}
\label{sec:discussion}

\subsection{Two fundamental failure modes}
\label{sec:failure-modes}

Our results identify two distinct mechanisms by which information is lost:

\textbf{1. Information destroyed (JPEG cascade / compaction loss)}

The single-pass experiment (\S\ref{sec:singlepass}) quantifies this directly: one compaction
pass renders the compacted zone near-dead (0--7\% recall). Keywords survive
in the summary (82--93\% grep recall) but are not retrievable. Multi-pass
strategies like Incremental compound this loss across cycles.

The strategy comparison (\S\ref{sec:strategy}) confirms this at scale: Brutal (S1) demonstrates total destruction
with recall dropping to 1.2\% at 5M tokens (61 compaction cycles). Each
cycle re-summarizes the entire history, amplifying loss exponentially.
Incremental (S2) fares better but remains heavily recency-biased --- at 1M
tokens, 94\% of its recall comes from the most recent quintile (Q5), with
the first four quintiles contributing near-zero.

The cross-model validation (\S\ref{sec:sonnet}) strengthens this finding: Sonnet 4.6, with
92.5\% baseline recall and no Lost-in-the-Middle effect, still drops to 21\% at
C4. If the loss were purely attentional, a more capable model should resist
proportionally better --- but it does not at severe compaction levels.

\textbf{2. Information preserved but ignored (attention dilution)}

The calibration's grep-LLM gap (86\% grep vs 67\% LLM at $\delta=0.42$) shows that facts
present in the context are not always found. The single-pass remaining-zone degradation
(73\% $\to$ 53\% as compaction increases) shows that injecting noise (re-padding)
dilutes attention even for intact facts.

The strategy comparison (\S\ref{sec:strategy}) reveals that attention dilution operates at the strategy level too.
Frozen (S3) preserves early facts as immutable summaries --- yet recall still
drops from 30\% at 500K to 10\% at 5M. The information is structurally intact
(frozen summaries are never re-compressed), but the model struggles to locate
relevant facts among 88 frozen blocks competing for attention with 400+
recent messages.

These failure modes are orthogonal and, critically, \textbf{both worsen with
scale}:
\begin{itemize}
  \item Frozen strategies prevent destruction but suffer attention dilution
    (too many summaries to scan)
  \item Incremental strategies destroy information but keep the context clean
    (fewer blocks to attend to)
  \item The optimal strategy must balance both failure modes --- and our strategy
    comparison shows that neither Frozen nor FrozenRanked achieves this
    balance at multi-megaToken scales
\end{itemize}

\subsection{The measurement problem}
\label{sec:measurement-problem}

The calibration findings (\S\ref{sec:calibration}) have implications beyond our own benchmark:

\textbf{Batch size sensitivity} means that compaction evaluations using different
values of $Q$ are not comparable. Worse, the direction of the Q-effect depends
on context quality (\S\ref{sec:q-inversion}): $Q=10$ outperforms $Q=1$ in rich contexts, but $Q=1$
outperforms $Q=5$ in severely compacted contexts. A strategy showing ``80\%
recall'' at $Q=10$ and another showing ``60\% recall'' at $Q=1$ may be equally
effective --- or the ranking may invert entirely.

\textbf{Category hierarchy} means that the mix of easy vs hard facts determines
the baseline. Testing compaction only on ``single-session-user'' facts (80--100\%
baseline) will show modest degradation. Testing on ``knowledge-update'' facts
(60--65\% baseline) will show severe degradation from a lower starting point.

\textbf{The grep-LLM gap} provides a free diagnostic: if grep finds a fact but the
LLM doesn't recall it, the problem is attention, not information loss. This
distinction is critical for choosing mitigation strategies (restructure
context vs improve compaction).

\subsection{Why global recall is misleading}
\label{sec:global-misleading}

A global recall score averages over spatial positions and fact categories. Our
calibration results show this hides structural information:
\begin{itemize}
  \item Categories range from 100\% (single-session-user) to 20\% (preferences)
  \item Spatial position matters: middle facts are harder to recall
  \item Batch size shifts the curve by up to 20pp, and the direction reverses under compaction
\end{itemize}

The single-pass experiment amplifies this: compacted-zone recall (0--7\%) and remaining-zone recall
(53--73\%) tell very different stories that a global 40\% would hide.

\textbf{Recommendation}: Always report spatially-resolved and category-resolved
metrics alongside global scores.

\subsection{Recall source: user queries vs chain of thought}
\label{sec:cot}

Our protocol tests recall via explicit user questions (``What was the server
IP?''). In real usage, context recall is predominantly driven by the model's
\textbf{chain of thought} --- during reasoning, the model implicitly scans the
context for relevant information. The access pattern likely differs:
\begin{itemize}
  \item Explicit queries target a single fact at a time
  \item Chain-of-thought may access 2--3 facts per reasoning step
  \item The type of information sought varies (code context vs decisions vs config)
\end{itemize}

Our calibration (\S\ref{sec:calibration}) and single-pass (\S\ref{sec:singlepass}) experiments cover $Q=1$ and $Q=5$
(and $Q=10$ for \S\ref{sec:calibration}); the strategy comparison (\S\ref{sec:strategy}) uses $Q=5$ throughout. This
spans the realistic range for explicit queries, but neither $Q$ approximates
the implicit, multi-fact access pattern of chain-of-thought. Instrumenting
a real agent to observe context access patterns (coding sessions,
conversational use, RAG-augmented workflows) is a promising direction for
future work.

\subsection{The limits of architectural optimization}
\label{sec:arch-limits}

The strategy comparison shows that FrozenRanked (S4) consistently
outperforms Frozen (S3) in the mean (S4--S3 gap of +13.2pp at 500K, +6.6pp
at 2M, +3.3pp at 3.5M, +1.6pp at 5M; the gap shrinks to +0.5pp at 1M
where one S3 replicate happens to recall unusually well). The advantage
exists but is modest at large scale, and the standard deviations (cf.\
\Cref{fig:phase-d-final}) overlap at every checkpoint past 500K --- the strict ordering
S4 $>$ S3 holds in the means but is not always significant on n=4--6
replicates.

The reason the gap stays modest is structural: at 5M tokens ($\sim$88 compaction
cycles), the frozen summaries together occupy only a small fraction of the
context --- well below the merge budget threshold. \textbf{Merges are rarely
triggered}, so S3 and S4 are functionally close. The merge strategies
that genuinely differentiate FrozenRanked from Frozen only activate when
frozen summaries exceed their budget, which requires conversations
exceeding $\sim$10M tokens ($\sim$168 cycles).

This suggests that \textbf{the bottleneck is not compression quality but attention
capacity}. Whether frozen summaries are merged sequentially or
hierarchically matters little if the model cannot effectively attend to
dozens of summary blocks in the first place. The attention dilution problem
identified in the single-pass experiment (\S\ref{sec:attention-dilution}) operates at the architectural level: more
preserved summaries means more context to scan, which degrades retrieval
for all facts --- frozen or not.

The implication is clear: improving compaction \textit{architecture} (how summaries
are organized and merged) yields diminishing returns. The next step requires
improving compaction \textit{interface} --- how preserved information is structured
for efficient retrieval (see \S\ref{sec:future}).

\subsection{Predictive modelling: what drives recall?}
\label{sec:predictive}

A logistic regression fitted across the \S\ref{sec:calibration} calibration and \S\ref{sec:singlepass} single-pass
datasets (n=3{,}852, AUC=0.84) confirms the expected drivers in additive
log-odds form: questions per prompt $Q$ (+0.086 / question), compaction
percent ($-0.054$ / percent), a positive compaction $\times$ position interaction
(later facts less damaged), and U-shaped position effect (Lost-in-the-Middle).
Density is not independently significant once $Q$ and position are controlled.
A gradient-boosted comparison on the same data scores AUC=0.64, so no
hidden non-linearities are missed by the additive model.

The model's main practical use is to express recall in compacted contexts
as simple additive rules of thumb rather than a black box. Full
specification, coefficients with confidence intervals, and the XGBoost
comparison are in Appendix B. \textbf{Caveat}: the 4{,}812 observations are not
mutually independent (same context, same facts repeated under different
conditions), so the reported p-values should be read as descriptive
indicators rather than rigorous significance tests; a mixed-effects
re-fit with a random intercept per fact is left for future work.

\subsection{Implications for production systems}
\label{sec:production}

Most production AI assistants (Claude Code, Cursor, Windsurf) use single-shot
or incremental compaction. Our results suggest:

\begin{enumerate}
  \item \textbf{Even light compaction is costly}: 5\% compaction costs 6--7pp recall.
    Systems should minimize compaction frequency.
  \item \textbf{The compacted zone is lost}: Don't expect the summary to be queryable
    for specific facts. If precise recall matters, extract facts to RAG
    \textit{before} compacting.
  \item \textbf{Re-padding noise hurts}: After compaction frees space, what fills that
    space matters. Random conversation is noise; structured information would
    be better.
  \item \textbf{Frozen summaries help but don't scale}: Freezing summaries preserves
    early facts (substantial improvement over Incremental at small scales, \S\ref{sec:strategy-results}) but the
    advantage narrows at scale as attention dilution dominates.
  \item \textbf{Batch size in evaluation matters}: Internal benchmarks should test at
    multiple values of $Q$ to avoid misleading conclusions.
  \item \textbf{Grep is not a valid proxy for recall}: Keywords surviving in summaries
    (82--93\%) vastly overestimates actual recoverability (0--7\%). Production
    systems should not rely on keyword presence as a quality signal.
\end{enumerate}
